\section{Introdução}
\label{sec:introducao}

% Sugestão de sequência: contexto -> problema -> evidência -> lacuna -> pergunta.
Contextualize a transformação ou o fenômeno investigado e delimite seu recorte.
Explique por que o problema importa para a mobilidade aérea urbana e sustente a
motivação com fontes verificáveis. Por exemplo, uma revisão pode ser citada
entre parênteses \citep{garrow2021uam}.

Apresente o que a literatura já explica e identifique precisamente o que ainda
não foi respondido. Encerre a seção com o objetivo e a pergunta de pesquisa.

\begin{quote}
\textbf{Pergunta de pesquisa:} como [fator] influencia [resultado] no contexto
de [recorte]?
\end{quote}

% Remova este comentário depois de verificar:
% - o recorte está explícito;
% - a lacuna decorre das referências apresentadas;
% - a pergunta pode ser respondida pelo método proposto.
